% Build the layer, Chapter 2, third build: revoking a card.
% Inserted after the second build (\S\ref{sec:anchor-build-holder}), whose
% gate it builds on. The design it builds part of is \S\ref{sec:revocation}'s.
%
% No code or output in this file is typed by hand. \pdexcerpt, \pdrun and
% \pdrunlines print files under excerpts/, which the kernel's
% tools/excerpts.mjs cuts from its source and its recorded runs (see
% figures/pd-listing.tex). To change a listing, change the kernel and cut
% again. The prose names functions, never line numbers, so that a listing
% which moves in its file does not make a sentence wrong.
%
% \pdlistingroom{n} before a listing of n lines keeps it on one page. When a
% listing changes length, change the number with it.
\subsection{Build the layer: revoking a card}
\label{sec:anchor-build-taken-back}

\noindent\emph{Express lane: a gate keeps a list of the cards that were
revoked. Looking up the card it is presented is not enough, because the
cards delegated from that card have ids of their own. The gate that works
looks up every card of a chain. It may prune an entry from its list once the
entry's card has expired, and not before.}

\S\ref{sec:revocation} set out a design for revoking a card, and it
has many parts. This section builds the part that one gate needs, in
\pdkernelname{}: a revocation list, and the two rules for it. One says which
cards a gate looks up. The other says which entries may be pruned. The
filter is not built, and the gossip that carries a revocation to other gates
is the next section's (\S\ref{sec:anchor-build-told}).

\subsubsection{Naive: look up the card that is presented}

The gate of \S\ref{sec:anchor-build-holder} checks a chain's cards and
makes their holder answer a challenge. It cannot refuse a card that is
valid. So give it a list of the cards that were revoked, and have it
look up the card it is presented. A card is listed by its id, which
\S\ref{sec:anchor-build} gave it: a serial number for the harbor's own
card, and the hash of the signature for a delegated card.

\pdlistingroom{21}
\pdexcerpt[caption={The first attempt. The card it looks up is the last of
the chain.},label={alg:anchor-taken-last}]{ch02c-last}

\texttt{revoke} adds the id of the last card of a chain to the list, and
\texttt{admit} refuses a chain whose last card is listed. Everything
else is left to the gate that this one wraps.

\subsubsection{Break: the cards delegated from a revoked card}

Heron's private key has leaked. \S\ref{sec:anchor-build-holder} ended
by saying where the key is, in his memory, and Wren has read it there
along with his cards. Alice
learns of it and asks the gate to revoke the card she delegated to him.
Finch holds a card that Heron delegated some time before.

\pdlistingroom{27}
\pdrun[caption={A card is revoked, and the cards delegated from it are
still accepted.},label={alg:anchor-taken-break}]{revoked-last-card}

The gate does refuse Heron's card, even when it is presented by someone
with his key. Then it admits Finch, because the card Finch presents is his
own, with an id that was never on the list. It admits Wren for the same
reason, on a card that did not exist when Alice went to the gate. With
Heron's key she can sign as many as she likes.

Revoking a card by its id covers that one card. The rights it granted are
still usable, through every card delegated from it.
Listing~\ref{alg:anchor-taken-links}, further on, revokes every set of
cards in a chain of three. At this gate the chain is still admitted for
three of the seven sets. \emph{Now you try:}
Exercise~\ref{ex:anchor-taken-count} asks for that count by hand.

\subsubsection{Fix: every card of the chain}

The check moves into the decision core, and it runs over every card.

\pdlistingroom{22}
\pdexcerpt[caption={The core's check that no card of a chain was
revoked.},label={alg:anchor-taken-standing}]{ch02c-not-revoked}

\texttt{checked\_cards} gives the id of each card in the chain with its
number, the harbor's card first. The core refuses the chain for the first
card that was revoked, and names it.

The core keeps no list, so it cannot look anything up. The gate does
that, and passes the core two lists: the ids it looked up, and those of
them it found revoked. The ids the core compares them with are its own. It
computed each one when it verified that card's signature, and it refuses a
presentation if the gate did not look up every one. A gate that looks up the
last card alone is refused. So is a gate that computes a card's id in some
other way.

In Listing~\ref{alg:anchor-taken-gate}, \texttt{admit} is the gate of
Listing~\ref{alg:anchor-holder-gate} with three more lines. It computes the
id of every card it is presented, and asks its list which of them are
revoked.

\pdlistingroom{35}
\pdexcerpt[caption={The gate. Every card is looked up, in the same turn as
the decision.},label={alg:anchor-taken-gate}]{ch02c-gate}

\texttt{revoke} passes the chain to the core, where the first attempt read
the last card and trusted it. The reason is where a card's fields are. Its
id is the hash of its
signature, and its expiry is in its claims. Text with the real
card's signature over other claims has the real card's id, and any expiry
its writer likes. Only the signature check tells the two apart. So the
core checks the chain, and returns the entry to add: the card's
id and the card's own expiry. The next break shows why the expiry matters.

% A page may break here. See the note in anchor-build-the-verifier.tex.
\pagebreak[0]
\keyidea[Revoking a card revokes every chain through it]{A card is a link
in every chain delegated from it. A gate that admits a chain accepts each
link of it, so it has to look up each one. Finch did nothing wrong, and his
card stopped working with Heron's.}

\subsubsection{Break again: a list pruned by age}

The list cannot grow for ever. A harbor that runs for a year revokes
a great many cards, and nearly all of them expired long ago. A cache
suggests a rule for that: remove an entry once it is old.

\pdlistingroom{11}
\pdexcerpt[caption={A list pruned by age.},label={alg:anchor-taken-aged}]{ch02c-ttl}

Alice revokes Heron's card three minutes into its ten, as she did at
the start of this section.

\pdlistingroom{20}
\pdrun[caption={A revoked card is accepted
again.},label={alg:anchor-taken-forgotten}]{pruned-ttl}

At $t=500000$, eight minutes and twenty seconds in, the entry is over
five minutes old and the list removes it. Heron's card has a hundred
seconds left, and the gate accepts it. Nothing was forged. The gate pruned
an entry it still needed.

\subsubsection{Fix: an entry is kept while its card is valid}

An entry stops a card being accepted. Once the card has expired, the gate
refuses it as expired, and the entry is no longer needed. Only then may it
be removed.

\pdlistingroom{17}
\pdexcerpt[caption={When an entry may be pruned.},label={alg:anchor-taken-spent}]{ch02c-prunable}

\pdlistingroom{14}
\pdexcerpt[caption={The list asks the core about each
entry.},label={alg:anchor-taken-forget}]{ch02c-prune}

\texttt{prunable} is why an entry records the expiry of its card. The list
passes each entry to the core, and removes the ones the core allows.

Design invariant~\ref{def:revocation-truncation} gives a reason for this
rule: no card outlives its parent, so nothing delegated from a pruned
entry's card can still be valid. This code needs less. It checks
the expiry of every card in a chain, so a chain through a revoked card is
refused once that card has expired, whatever the cards delegated from
it claim. A verifier that checked only the last card's expiry would need the
stronger reason.

\subsubsection{Verify: the same runs, then every case}

The leaked key first (Listing~\ref{alg:anchor-taken-fixed}).

\pdlistingroom{16}
\pdrunlines[caption={The leaked key, against the fix. Every chain through
Heron's card is refused.},label={alg:anchor-taken-fixed}]{revoked-core}{1,13-27}{1 and 13\textendash 27}

Each refusal names card 2, which is Heron's, though the last card presented
is someone else's in two of the three.

Then the list pruned by age (Listing~\ref{alg:anchor-taken-kept}).

\pdlistingroom{13}
\pdrunlines[caption={The same run, against the fix. The entry is pruned
when the card has expired.},label={alg:anchor-taken-kept}]{pruned-core}{1,9-20}{1 and 9\textendash 20}

That run tries one pair of times. Listing~\ref{alg:anchor-taken-moments}
tries every pair of whole minutes in the card's ten: the minute it is
revoked, and the minute it is presented.

\pdlistingroom{22}
\pdrun[caption={Every pair of minutes, at the list pruned by
age.},label={alg:anchor-taken-moments}]{timings-ttl}

Each \texttt{!} is a card accepted after it was revoked. They are where
the presentation is more than five minutes after the revocation, which is
15 of the 66 pairs. Under the fix the count is none
(Listing~\ref{alg:anchor-taken-moments-core}).

\pdlistingroom{2}
\pdrunlines[caption={Every pair of minutes, against the
fix.},label={alg:anchor-taken-moments-core}]{timings-core}{1,19}{1 and 19}

The first break is checked in every case too.
Listing~\ref{alg:anchor-taken-links} revokes every set of cards in a chain of
three. The first attempt refuses the chain when the third card is among those
revoked, and at no other time. The core refuses it for all seven sets.

\pdlistingroom{19}
\pdrun[caption={Every set of cards of a chain of three that can be
revoked.},label={alg:anchor-taken-links}]{subsets}

Last, the guards. Build the core so that an entry may be pruned at any
time, and every presentation is admitted
(Listing~\ref{alg:anchor-taken-moments-mutant}).

\pdlistingroom{2}
\pdrunlines[caption={The fix with an entry allowed to be pruned at any
time.},label={alg:anchor-taken-moments-mutant}]{timings-prune-early}{1,19}{1 and 19}

Build it so that it checks the last card alone, and the count of
Listing~\ref{alg:anchor-taken-links} returns
(Listing~\ref{alg:anchor-taken-links-mutant}).

A third build of the core does not check which cards the gate looked up,
and the kernel names the test that fails on it. Two guards from
\S\ref{sec:anchor-build} are tried here as well. With the check of a
card's expiry removed, a revoked card is accepted again when its entry is
pruned, and a test fails. With the comparison of expiry times removed, none
does: the pruning rule does not depend on it.

\pdlistingroom{13}
\pdrunlines[caption={The fix with its check narrowed to the last
card.},label={alg:anchor-taken-links-mutant}]{subsets-last-card-only}{1,6-17}{1 and 6\textendash 17}

The earlier builds each had a ProVerif model. This one has a model, and
no proof (Listing~\ref{alg:anchor-taken-models}).

\pdlistingroom{3}
\pdrun[caption={ProVerif~2.05 on a gate that revokes one card and then
looks it up.},label={alg:anchor-taken-models}]{models-revocation}

The model, \texttt{revoked.pv}, is as small as the question allows. A
gate adds one card to its list and then looks that card up, so the card can
never be admitted. ProVerif does not prove that. It reasons with facts that
stay true once they are true, and ``this card is not in the list'' is
not such a fact. So it treats the \texttt{else} branch of a lookup as one
that may always be taken, and its answer in
Listing~\ref{alg:anchor-taken-models} is that the query ``cannot be
proved''. It reports no attack either: it found no run in which the card
is admitted.
Whether a revoked card stays refused depends on something being
absent, and ProVerif~2.05 does not reason about absence.
\pdscope{Listings~\ref{alg:anchor-taken-links}
and~\ref{alg:anchor-taken-moments} cover a chain of three cards and a card
valid for ten minutes, at one gate with one clock. They are exhaustive over
those small cases, and they do not prove that larger ones behave the same.}

\begin{pdboundary}
\setlength{\parskip}{4pt plus 1pt}
The gate does not check who is revoking a card. Anyone who can present a
chain can have its last card revoked. Wren, with copies of Heron's
cards and none of his keys, could not use them after the last build. Here
she can revoke his card and lock him out
(Exercise~\ref{ex:anchor-taken-who}), until
\S\ref{sec:anchor-build-warranted} checks who is revoking. The core also
trusts its gate. It
checks that the gate claims to have looked up every card, and it cannot
check that the list answered truthfully.

This is one gate. A second gate that has not been informed of a revocation
goes on accepting the card, as Figure~\ref{fig:anchor-revocation-gossip}
shows. Spreading the revocation is \S\ref{sec:anchor-build-told}. The
pruning rule uses one clock, and two gates whose clocks differ need the
margin that Design invariant~\ref{def:revocation-truncation} names. The
list has no size limit of its own: it holds every revoked card that has not
yet expired.
\end{pdboundary}
